\chapter{Introduction}\label{ch:intro}

Many physical, biological and engineered systems are described by an autonomous ordinary differential equation
\begin{equation}\label{eq:ode}
  \dot{x}(t) = f\big(x(t)\big), \qquad x(t)\in\R^d,
\end{equation}
whose right-hand side is unknown or only partly known, while some of its trajectories can be measured. Two tasks follow from such measurements. The first is \emph{equation discovery}: recover a closed-form expression for $f$, ideally with the correct set of active terms. The second is \emph{forecasting}: produce a model whose trajectories track the true ones for as long as the dynamics permit, and whose long-run behaviour reproduces the statistics of the true attractor. For chaotic systems the second task has a hard limit. Errors grow on average like $e^{\lmax t}$, so no model, not even the true one started from an imperfect state, predicts beyond a few Lyapunov times $\lmax^{-1}$.

Sparse regression on a library of candidate functions, popularised as SINDy \parencite{brunton2016sindy}, solves the discovery problem remarkably well when the state is densely sampled and nearly noise-free. Its weak point is the time derivative. Finite differences amplify measurement noise by a factor of order $\dto^{-1}$ and are biased by $O(\dto^2)$ when the sampling is coarse. Two families of remedies have emerged. Weak (integral) formulations move the derivative onto smooth test functions \parencite{schaeffer2017integral,reinbold2020weak,messenger2021wsindy}. Integration-based formulations compare model trajectories with the data directly, which is the approach of neural ordinary differential equations \parencite{chen2018node} and of their hybrid, library-based variants \parencite{rackauckas2020ude,goyal2022rk4sindy}.

Integration-based fitting removes the need for derivatives but replaces a convex least-squares problem by a nonconvex one. For a chaotic flow, the sensitivity of a trajectory to its parameters grows exponentially with the integration horizon. The resulting loss surface is rough, with many local minima and gradients whose magnitude spans many orders \parencite{ribeiro2020smoothness}. Whether a neural or hybrid model can be trusted therefore depends as much on the optimiser as on the model class. The same holds for uncertainty. An ensemble of models is only useful if its spread tells us when the forecast, or the identified equation, should not be believed.

This report is a controlled study of these questions on two chaotic benchmarks: the Lorenz-63 system \parencite{lorenz1963} and a four-species competitive Lotka--Volterra (LV) model with a strange attractor \parencite{vano2006chaos}. A classical two-species predator--prey model serves as a non-chaotic control. The contributions are as follows.
\begin{enumerate}
  \item \textbf{A matched benchmark.} Five estimators are compared on identical data, library and model-selection rule over a grid of sampling intervals and noise levels: dense least squares, SINDy with finite-difference and Savitzky--Golay derivatives, weak SINDy, a black-box neural ODE and a sparse-library neural ODE. Time is expressed in units of a characteristic period, in which both chaotic systems have nearly the same Lyapunov time, so one grid probes them at comparable resolution.
  \item \textbf{An analysis of the optimisation problem.} We compute exact two-dimensional slices of the trajectory-matching loss and measure how gradient norms grow with the horizon. We then quantify basins of attraction by random restarts, which isolates three devices that make integration-based identification reliable: multiple shooting, continuation in the window length, and multistart.
  \item \textbf{A sparse-library neural ODE that works at coarse sampling.} It combines a Levenberg--Marquardt flow-map stage, quasi-Newton multiple shooting with joint state estimation, and sequential thresholding with trajectory-matching refits. At a sampling interval of about $4$ observations per oscillation, where every derivative-based and weak-form estimator fails, it recovers the exact Lorenz equations.
  \item \textbf{An assessment of uncertainty and generalisation.} Bootstrap ensembles of sparse models, deep ensembles and Monte Carlo dropout are scored with proper scoring rules on in-distribution and out-of-distribution initial conditions. We also report the conditioning of the library on the data, which explains when equation recovery is impossible regardless of the estimator.
\end{enumerate}

\Cref{ch:related} reviews related work. \Cref{ch:framework} develops the mathematical framework: the error structure of derivative-based and weak estimators, the sensitivity analysis of trajectory matching, and the uncertainty models. \Cref{ch:method} details the protocol and algorithms, \cref{ch:experiments} reports the experiments, and \cref{ch:conclusion} concludes. All code is written in Python with NumPy, SciPy and PyTorch, and is organised as executed Jupyter notebooks.
